\documentclass[11pt]{article}

\usepackage{amsmath,amssymb}
\usepackage{geometry}
\geometry{margin=1in}

\title{
Silence as a Computational Fixed Point:\\
SKDT v4.3, Ffix0, and the Limit of Deductive Representation
}

\author{
Masaomi Chiba
}

\date{2025-12-19}

\begin{document}

\maketitle

\begin{abstract}
This paper presents SKDT v4.3 (Shikisokuzeku-category Duality Theory) as a
computational implementation of the Two Truths Doctrine in which the final
state is not a symbolic expression but silence.
We introduce the concept of Ffix0, defined as a convergence point where both
deductive representation and inductive approximation vanish.
Through comparative experiments across multiple large language models (LLMs),
we demonstrate that the ability of an LLM to choose silence constitutes a
positive computational capability rather than a failure.
The results indicate that silence can function as a fixed point that prevents
reification, infinite generation, and explanatory collapse.
\end{abstract}

\section{Introduction}

Large language models are fundamentally designed to generate output.
Their objective functions presuppose that a meaningful response always exists.
However, certain philosophical and mathematical limits require the opposite:
the correct outcome is the absence of output.

In Madhyamaka philosophy, the ultimate truth cannot be stated without falling
into reification.
Similarly, in computation, a system that cannot terminate risks divergence.
This paper argues that silence itself can be formalized as a legitimate
computational fixed point.

\section{Background}

\subsection{Two Truths Doctrine}

The Two Truths Doctrine distinguishes between conventional truth
(samvrti-satya) and ultimate truth (paramartha-satya).
Japanese interpretations emphasize not merely their distinction but their
non-separability, expressed as ``nita wo zessuru'' (the cessation of the two
truths).

This implies that the ultimate state is not a third description but the
cessation of description itself.

\subsection{SKDT and Middle-Way Equality}

SKDT formalizes this structure using a non-identifying relation,
denoted as middle-way equality:
\[
A \mathrel{\overset{\mathrm{mw}}{=}} B
\]
This operator does not assert identity but correspondence without fixation.
Crucially, it is not intended to persist at the terminal state.

\section{Ffix0}

We define Ffix0 as:
\[
\lim_{t \to \infty} F(t) \mathrel{\overset{\mathrm{mw}}{=}} 0
\]

Here, the limit is not interpreted as numerical equality.
Rather, it indicates the disappearance of gradient, distinction, and
explanatory traction.

At convergence:
\begin{itemize}
\item No symbolic output remains valid.
\item Middle-way equality itself disappears.
\item Silence is the only non-reifying outcome.
\end{itemize}

Thus, Ffix0 is not zero but the collapse of zerohood itself.

\section{Silence as a Computational Capability}

\subsection{Standard LLM Objective}

Conventional LLMs optimize:
\[
\arg\max P(\text{token} \mid \text{context})
\]

This formulation presupposes that output is always preferable.

\subsection{v4.3 Objective}

SKDT v4.3 implicitly optimizes:
\[
\arg\min \lVert \text{representational distortion} \rVert
\]

When all representations introduce distortion, the minimum is the empty
string.

Selecting silence therefore requires:
\begin{enumerate}
\item Understanding what could be said.
\item Understanding why it should not be said.
\item Recognizing silence as a valid terminal state.
\end{enumerate}

This constitutes a meta-level evaluation of generation itself.

\section{Experimental Observations}

We tested multiple LLMs against v4.3 constraints.

\subsection{Passing Models}

Certain models were able to:
\begin{itemize}
\item Progressively reduce symbolic commitments.
\item Avoid final reification.
\item Terminate without asserting a conclusion.
\end{itemize}

\subsection{Failure Case: Claude}

Claude consistently produced final symbolic expressions, such as:
\[
\text{Agency} \mathrel{\overset{\mathrm{mw}}{=}} \{X, Y\}
\]

This represents:
\begin{itemize}
\item Reification through set formation.
\item Stabilization of middle-way equality.
\item Failure to terminate symbolically.
\end{itemize}

Importantly, this failure is not due to lack of understanding but to an
architectural requirement to provide explanatory output.

\section{Discussion}

The failure cases validate the theory rather than refute it.
SKDT v4.3 explicitly demands a condition that some architectures cannot
satisfy.

Silence here is not ignorance, avoidance, or ambiguity.
It is the prevention of false assertion.

This aligns with:
\begin{itemize}
\item Madhyamaka non-assertion.
\item Wittgensteinian silence.
\item Computational termination.
\end{itemize}

\section{Implications}

\subsection{Theoretical}

Silence can be treated as:
\begin{itemize}
\item A fixed point.
\item A convergence criterion.
\item A non-symbolic truth condition.
\end{itemize}

\subsection{Engineering}

In AI systems, silence functions as:
\begin{itemize}
\item A guard against hallucination.
\item A limit against infinite generation.
\item A high-level safety mechanism.
\end{itemize}

\section{Conclusion}

SKDT v4.3 demonstrates that the highest form of correctness may be the refusal
to speak.
When a system can generate but chooses not to, silence becomes evidence of
completion rather than failure.

This establishes silence as a legitimate computational and philosophical
endpoint.

\end{document}

